% ============================================================
% Rapport mini-projet - Agent-Assisted Data Warehouse Design
%                       and ETL Automation
% Auteur : HALIMI Mohamed Salah Eddine
% Encadreur : Cpt. HAMOUDA Sidehoum
% Specialite : SIAD (Systeme d'Information et Aide a la Decision)
% ============================================================
\documentclass[12pt, a4paper]{report}

% ── Encodage et langue ──────────────────────────────────────────
\usepackage[utf8]{inputenc}
\usepackage[T1]{fontenc}
\usepackage[french]{babel}
\usepackage{lmodern}

% ── Mise en page ────────────────────────────────────────────────
\usepackage[a4paper, top=2.5cm, bottom=2.5cm, left=3cm, right=2.5cm]{geometry}
\usepackage{setspace}
\onehalfspacing
\usepackage{fancyhdr}
\pagestyle{fancy}
\fancyhf{}
\fancyhead[L]{\textsc{EMP}}
\fancyhead[R]{\leftmark}
\fancyfoot[C]{\thepage}
\renewcommand{\headrulewidth}{0.4pt}

% ── Graphiques et tableaux ──────────────────────────────────────
\usepackage{graphicx}
\usepackage{float}
\usepackage{tikz}
\usetikzlibrary{shapes.geometric, arrows.meta, positioning, fit, backgrounds, shadows}
\usepackage{pgfplots}
\pgfplotsset{compat=1.18}
\usepackage{booktabs}
\usepackage{tabularx}
\usepackage{longtable}
\usepackage{array}
\usepackage{multirow}
\usepackage{caption}

% ── Code et listings ────────────────────────────────────────────
\usepackage{listings}
\usepackage{xcolor}
\definecolor{codebg}{RGB}{248,250,252}
\definecolor{codekw}{RGB}{99,102,241}
\definecolor{codestr}{RGB}{5,150,105}
\definecolor{codecom}{RGB}{100,116,139}
\lstset{
  backgroundcolor=\color{codebg},
  basicstyle=\ttfamily\footnotesize,
  keywordstyle=\color{codekw}\bfseries,
  stringstyle=\color{codestr},
  commentstyle=\color{codecom}\itshape,
  numbers=left, numberstyle=\tiny\color{gray},
  breaklines=true, frame=single, captionpos=b, tabsize=2,
}

% ── Liens et references ─────────────────────────────────────────
\usepackage[colorlinks=true, linkcolor=black, urlcolor=blue, citecolor=blue]{hyperref}
\usepackage{cleveref}
\usepackage{enumitem}

% ── Titres de chapitres ─────────────────────────────────────────
\usepackage{titlesec}
\titleformat{\chapter}[display]
  {\normalfont\Huge\bfseries}{\chaptertitlename\ \thechapter}{20pt}{\Huge}

% ── Acronymes ───────────────────────────────────────────────────
\usepackage{acronym}

% ============================================================
\begin{document}

% ────────────────────────────────────────────────────────────
% PAGE DE GARDE
% ────────────────────────────────────────────────────────────
\begin{titlepage}
  \centering
  \vspace*{0.5cm}
  {\large Ministère de la Défense Nationale}\\[0.3cm]
  {\large État-Major de l'Armée Nationale Populaire}\\[0.3cm]
  {\large École Militaire Polytechnique}\\[0.2cm]
  {\large \textit{Chahid Abderrahmane Taleb}}\\[1.5cm]

  \rule{\linewidth}{0.5mm}\\[0.4cm]
  {\Huge \bfseries Rapport du mini-projet}\\[0.4cm]
  \rule{\linewidth}{0.5mm}\\[1.2cm]

  {\LARGE \bfseries Agent-Assisted Data Warehouse Design\\[0.2cm] and ETL Automation}\\[0.5cm]

  \vspace{1cm}
  {\large \textit{Conception assistée par agents IA et automatisation\\
  de la chaîne ETL pour entrepôts de données}}\\[1.5cm]

  \begin{minipage}{0.45\textwidth}
    \begin{flushleft}\large
      \textbf{RÉALISÉ PAR :}\\[0.3cm]
      E.O.C : HALIMI Mohamed Salah Eddine
    \end{flushleft}
  \end{minipage}
  \hfill
  \begin{minipage}{0.45\textwidth}
    \begin{flushright}\large
      \textbf{ENCADRÉ PAR :}\\[0.3cm]
      Cpt. HAMOUDA Sidehoum
    \end{flushright}
  \end{minipage}

  \vfill
  {\large \textbf{Spécialité :} Génie informatique}\\[0.2cm]
  {\large \textbf{Option :} SIAD (Système d'Information et Aide à la Décision)}\\[1cm]
  {\large Année universitaire : 2025/2026}
\end{titlepage}

\cleardoublepage

% ────────────────────────────────────────────────────────────
% DEDICACES
% ────────────────────────────────────────────────────────────
\chapter*{Dédicaces}
\addcontentsline{toc}{chapter}{Dédicaces}
\thispagestyle{empty}

\begin{flushright}\itshape

À mes chers parents,\\
qui m'ont toujours soutenu par leurs sacrifices, leur amour et leurs prières.\\
Tout ce que je suis aujourd'hui, je le leur dois.\\[0.5cm]

À mes frères et sœurs,\\
pour leur affection et leur soutien constant tout au long de mon parcours.\\[0.5cm]

À mon encadreur, le Capitaine HAMOUDA Sidehoum,\\
dont l'expertise, la rigueur scientifique et la disponibilité ont été\\
des sources précieuses d'apprentissage et d'inspiration.\\
Je lui exprime ma plus profonde gratitude pour la confiance\\
qu'il m'a accordée et pour la qualité de son encadrement.\\[0.5cm]

À tous les enseignants de l'École Militaire Polytechnique,\\
pour la qualité de la formation qu'ils m'ont transmise.\\[0.5cm]

À mes camarades de promotion,\\
pour les moments partagés, l'entraide et la solidarité.\\[0.5cm]

À tous ceux qui, de près ou de loin, ont contribué\\
à la réussite de ce travail.\\[1cm]

\textbf{Halimi Mohamed Salah Eddine}
\end{flushright}

\cleardoublepage

% ────────────────────────────────────────────────────────────
% REMERCIEMENTS
% ────────────────────────────────────────────────────────────
\chapter*{Remerciements}
\addcontentsline{toc}{chapter}{Remerciements}
\thispagestyle{empty}

Avant tout, je rends grâce à \textbf{Dieu le Tout-Puissant} de m'avoir donné la santé,
la patience et la volonté nécessaires pour mener à bien ce modeste travail.

J'exprime ma profonde gratitude à mes \textbf{parents} pour leur soutien
inconditionnel, leurs sacrifices et leurs prières qui ont constitué le fondement
de ma réussite tout au long de mon parcours.

Mes plus sincères remerciements s'adressent à mon encadreur,
\textbf{le Capitaine HAMOUDA Sidehoum},
pour son encadrement de qualité, sa disponibilité, ses conseils avisés et la confiance
qu'il m'a témoignée. Sa rigueur scientifique, son expertise dans les domaines de
l'aide à la décision et des systèmes d'information ont été des éléments
déterminants dans l'aboutissement de ce projet.

Je tiens à exprimer ma reconnaissance aux \textbf{membres du jury} qui me font
l'honneur d'examiner et d'évaluer ce travail.

Je remercie également l'ensemble du corps enseignant de
l'\textbf{École Militaire Polytechnique}, et en particulier le département
SIAD, pour la qualité de la formation reçue durant ces années.

Enfin, je remercie chaleureusement mes \textbf{camarades de promotion} pour
leur aide, leur soutien constant et l'esprit de solidarité qui a régné tout
au long de cette formation.

\cleardoublepage

% ────────────────────────────────────────────────────────────
% RESUME / ABSTRACT
% ────────────────────────────────────────────────────────────
\chapter*{Résumé}
\addcontentsline{toc}{chapter}{Résumé}
\thispagestyle{empty}

La conception et l'industrialisation d'un \emph{Data Warehouse} demeurent
des activités complexes, chronophages et fortement dépendantes de l'expertise
humaine. Les projets d'entreposage de données impliquent un ensemble d'étapes
sensibles~: profilage des sources, modélisation dimensionnelle selon
la méthodologie de Kimball, génération du DDL T-SQL, écriture des pipelines
d'extraction-transformation-chargement (ETL), validation de la qualité des données
et restitution décisionnelle.

Ce mini-projet propose une plate-forme intelligente, intitulée
\textbf{Agent Data Warehouse}, qui automatise l'intégralité de cette chaîne
grâce à un système \emph{multi-agents} orchestré par \texttt{LangGraph}, intégrant
des grands modèles de langage (LLM) avec validation humaine en boucle
(\emph{Human-in-the-Loop}). L'architecture proposée combine~:
(i)~vingt-trois agents spécialisés~couvrant l'exploration des sources,
le calcul de la qualité, la modélisation en étoile, la critique automatique,
la génération T-SQL, l'exécution ETL avec auto-réparation, le suivi du lignage
et la génération d'un rapport décisionnel~;
(ii)~un assistant conversationnel \textbf{Atlas} fondé sur un mode \emph{patch
operations} et une cascade de fournisseurs LLM (GLM-5, OpenRouter, Ollama)
avec un \emph{smart parser} déterministe en repli~;
(iii)~une infrastructure de déploiement \emph{production-ready} (Docker
multi-étapes, reverse-proxy Caddy avec HTTPS automatique, observabilité
Prometheus / Grafana, tests automatisés à plusieurs niveaux).

Les expérimentations menées sur la base de données \textit{Northwind} démontrent
qu'un schéma en étoile complet, un DDL exécutable et un rapport Excel multi-feuilles
peuvent être produits sans intervention manuelle, tout en conservant un contrôle
explicite de l'utilisateur via la phase de validation humaine.

\textbf{Mots-clés~:} Data Warehouse, ETL, Kimball, agents IA, LangGraph,
LLM, Human-in-the-Loop, automatisation, OLAP.

\cleardoublepage

\chapter*{Abstract}
\addcontentsline{toc}{chapter}{Abstract}
\thispagestyle{empty}

Designing and industrializing a Data Warehouse remains a complex,
time-consuming task that heavily relies on human expertise. Such projects
involve multiple sensitive steps~: source profiling, Kimball dimensional
modeling, T-SQL DDL generation, ETL pipeline authoring, data-quality
validation and decision-oriented reporting.

This mini-project introduces an intelligent platform, \textbf{Agent Data
Warehouse}, that automates the entire chain through a multi-agent system
orchestrated by \texttt{LangGraph}, integrating Large Language Models (LLM)
with Human-in-the-Loop validation. The proposed architecture combines~:
(i) twenty-three specialized agents covering source exploration, data quality
scoring, star-schema modeling, automatic critique, T-SQL generation, ETL
execution with self-healing, lineage tracking and decision reporting~;
(ii) a conversational assistant \textbf{Atlas} based on a patch-operations mode
with a multi-provider LLM cascade (GLM-5, OpenRouter, Ollama) and a
deterministic smart parser as ultimate fallback~; (iii) a production-ready
deployment infrastructure (multi-stage Docker, automatic-HTTPS Caddy reverse
proxy, Prometheus/Grafana observability, multi-level automated tests).

Experiments on the Northwind database show that a complete star schema, an
executable DDL and a multi-sheet Excel report can be produced without manual
intervention while keeping explicit user control through the human-validation
step.

\textbf{Keywords~:} Data Warehouse, ETL, Kimball, AI agents, LangGraph,
LLM, Human-in-the-Loop, automation, OLAP.

\cleardoublepage

% ────────────────────────────────────────────────────────────
% TABLE DES MATIERES + LISTES
% ────────────────────────────────────────────────────────────
\tableofcontents
\cleardoublepage
\listoffigures
\addcontentsline{toc}{chapter}{Liste des figures}
\cleardoublepage
\listoftables
\addcontentsline{toc}{chapter}{Liste des tableaux}
\cleardoublepage

% ────────────────────────────────────────────────────────────
% LISTE DES ACRONYMES
% ────────────────────────────────────────────────────────────
\chapter*{Liste des acronymes}
\addcontentsline{toc}{chapter}{Liste des acronymes}
\begin{acronym}[OPENROUTER]
  \acro{API}{Application Programming Interface}
  \acro{BI}{Business Intelligence}
  \acro{CDC}{Change Data Capture}
  \acro{CSV}{Comma-Separated Values}
  \acro{DDL}{Data Definition Language}
  \acro{DML}{Data Manipulation Language}
  \acro{DQ}{Data Quality}
  \acro{DW}{Data Warehouse}
  \acro{ETL}{Extract-Transform-Load}
  \acro{ELT}{Extract-Load-Transform}
  \acro{FK}{Foreign Key}
  \acro{GLM}{General Language Model}
  \acro{HITL}{Human-in-the-Loop}
  \acro{HTTP}{HyperText Transfer Protocol}
  \acro{HTTPS}{HyperText Transfer Protocol Secure}
  \acro{JSON}{JavaScript Object Notation}
  \acro{JWT}{JSON Web Token}
  \acro{KPI}{Key Performance Indicator}
  \acro{LLM}{Large Language Model}
  \acro{OLAP}{Online Analytical Processing}
  \acro{OLTP}{Online Transactional Processing}
  \acro{PK}{Primary Key}
  \acro{REST}{Representational State Transfer}
  \acro{SCD}{Slowly Changing Dimension}
  \acro{SK}{Surrogate Key}
  \acro{SLA}{Service-Level Agreement}
  \acro{SQL}{Structured Query Language}
  \acro{SSE}{Server-Sent Events}
  \acro{TSQL}{Transact-SQL}
  \acro{UI}{User Interface}
  \acro{XAI}{eXplainable Artificial Intelligence}
  \acro{XML}{Extensible Markup Language}
\end{acronym}
\cleardoublepage

% ────────────────────────────────────────────────────────────
% INTRODUCTION GENERALE
% ────────────────────────────────────────────────────────────
\chapter*{Introduction générale}
\addcontentsline{toc}{chapter}{Introduction générale}
\markboth{Introduction générale}{}

\paragraph{Contexte.}
Dans un monde où les organisations produisent des volumes de données
sans précédent, la capacité à transformer ces flux hétérogènes en
\emph{connaissance actionnable} constitue un avantage concurrentiel décisif.
Le \ac{DW} demeure la pierre angulaire de la \ac{BI} d'entreprise~: il
agrège les données opérationnelles issues de multiples systèmes sources
(\ac{OLTP}), les nettoie, les modélise selon une logique dimensionnelle et
les met à disposition d'analystes, de tableaux de bord et de modèles
prédictifs.

\paragraph{Problématique.}
Concevoir un \ac{DW} de qualité reste cependant une tâche difficile.
Les ingénieurs doivent~:
\begin{enumerate}[label=(\alph*)]
  \item profiler des sources hétérogènes (SQL Server, MySQL, PostgreSQL,
        REST, fichiers CSV, sauvegardes \texttt{.bak})~;
  \item identifier la \emph{table de faits} et les dimensions selon la
        méthodologie de Kimball~;
  \item gérer les évolutions du schéma (\emph{schema drift}) et la
        \ac{SCD}~;
  \item écrire et maintenir des centaines de lignes de \ac{TSQL} pour les
        procédures ETL~;
  \item garantir la qualité, la traçabilité (lignage), la gouvernance,
        et produire des restitutions exploitables par les décideurs.
\end{enumerate}
Chacune de ces étapes mobilise des compétences spécialisées
(architecte data, développeur ETL, analyste BI), ce qui rallonge les délais
de mise en production et augmente le coût global.

\paragraph{Objectifs.}
Ce projet vise à concevoir et implémenter un système \emph{multi-agents}
qui automatise l'intégralité du cycle de vie d'un \ac{DW}, de l'ingestion
de la source brute jusqu'à la livraison du rapport décisionnel, tout en
maintenant l'humain en boucle aux étapes critiques (validation du modèle).
Plus précisément, les objectifs sont les suivants~:
\begin{itemize}
  \item proposer une architecture \emph{multi-agents} orchestrée par
        \texttt{LangGraph} couvrant les phases d'exploration, de qualité,
        de modélisation, de critique, de \ac{DDL}, d'ETL et de lignage~;
  \item concevoir un assistant conversationnel \textbf{Atlas} permettant
        à l'utilisateur de modifier le schéma en langage naturel, avec
        une cascade robuste de fournisseurs \ac{LLM}~;
  \item garantir un comportement déterministe en cas d'indisponibilité
        des modèles \ac{LLM} grâce à un \emph{smart parser} fondé sur des
        heuristiques métier (split de clés date, détection de mesures
        calculées, ajout de \ac{FK})~;
  \item fournir une infrastructure de déploiement \emph{production-ready}
        (Docker multi-étapes, \ac{HTTPS} automatique, observabilité,
        146 tests automatisés).
\end{itemize}

\paragraph{Contributions.}
Les principales contributions de ce travail sont~:
\begin{itemize}
  \item une architecture \emph{multi-agents} de 23 nœuds spécialisés
        avec \emph{checkpointing} et reprise après interruption~;
  \item un mécanisme de \emph{patch operations} permettant à un \ac{LLM}
        de produire des opérations atomiques de modification du modèle
        plutôt que de régénérer le schéma entier (10$\times$ moins de
        tokens, 100$\times$ plus fiable)~;
  \item un \emph{smart parser} déterministe capable d'inférer jusqu'à
        six opérations de modification à partir d'une demande Kimball
        complexe (multi-dates, mesures calculées, contraintes
        d'intégrité), sans aucun \ac{LLM}~;
  \item un rapport Excel décisionnel à dix feuilles intégrant des
        graphiques natifs et des agrégations multi-granularité~;
  \item un pipeline d'observabilité complet (Prometheus
        \texttt{/metrics}, dashboard Grafana, logs JSON, cache
        \ac{LLM}~LRU+TTL).
\end{itemize}

\paragraph{Plan du rapport.}
Le présent rapport est organisé en trois chapitres principaux~:
\begin{itemize}
  \item Le \textbf{Chapitre~\ref{chap:concepts}} introduit les concepts
        fondamentaux relatifs aux entrepôts de données, à la
        méthodologie de Kimball, aux pipelines \ac{ETL}, aux systèmes
        multi-agents et aux modèles de langage.
  \item Le \textbf{Chapitre~\ref{chap:etat}} dresse un état de l'art
        des outils et approches existants (dbt, Airbyte, Fivetran,
        Talend, agents \ac{LLM}) et positionne notre contribution.
  \item Le \textbf{Chapitre~\ref{chap:conception}} présente la
        conception détaillée du système \emph{Agent Data Warehouse},
        son architecture, ses composants et les résultats expérimentaux.
\end{itemize}

\cleardoublepage

% ============================================================
% CHAPITRE I : CONCEPTS FONDAMENTAUX
% ============================================================
\chapter{Concepts fondamentaux pour la conception assistée d'entrepôts de données}
\label{chap:concepts}

\section{Introduction}
Ce chapitre présente les concepts théoriques nécessaires à la
compréhension du projet. Il commence par positionner l'\ac{ETL} et
l'entreposage de données dans un contexte décisionnel, expose la
méthodologie de Kimball pour la modélisation dimensionnelle, puis
introduit les paradigmes de l'intelligence artificielle agentique
(\emph{multi-agents}, \texttt{LangGraph}, \ac{LLM}) qui sont au cœur
de la solution proposée.

\section{Entreposage de données et Business Intelligence}

\subsection{Définition et finalité d'un Data Warehouse}
Un \ac{DW} est, selon la définition canonique de Bill Inmon, une
collection de données~\textit{orientée sujet}, \textit{intégrée},
\textit{historisée} et \textit{non volatile}, conçue pour soutenir
le processus de décision. À la différence d'une base \ac{OLTP} dont la
vocation est de gérer des transactions à haute fréquence et faible
volume, un \ac{DW} est optimisé pour les requêtes analytiques massives
sur des plages temporelles étendues.

\subsection{Architecture en couches}
Une architecture \ac{DW} typique se compose de trois couches~:
\begin{itemize}
  \item \textbf{Couche source}~: bases \ac{OLTP}, fichiers plats,
        \acp{API} REST, sauvegardes binaires (\texttt{.bak}).
  \item \textbf{Couche d'intégration}~: pipelines \ac{ETL} qui
        extraient, nettoient, transforment et chargent les données.
  \item \textbf{Couche de présentation}~: tables dimensionnelles
        organisées en étoile ou en constellation, accessibles via
        des outils \ac{OLAP} et des tableaux de bord.
\end{itemize}

\subsection{Pipeline ETL versus ELT}
Historiquement, le paradigme \ac{ETL} consiste à transformer les
données avant leur chargement dans le \ac{DW}, ce qui demande une
\emph{staging area} intermédiaire. Avec l'avènement des entrepôts
massivement parallèles, l'approche \ac{ELT} a gagné du terrain~:
les données brutes sont chargées telles quelles dans l'entrepôt, puis
transformées \emph{in situ} via du \ac{SQL}. Notre projet adopte une
approche hybride~: extraction et transformation contrôlées par des
agents Python, chargement final en \ac{TSQL} natif sur SQL Server.

\section{Modélisation dimensionnelle de Kimball}
\label{sec:kimball}

\subsection{Étoile, flocon, constellation}
La méthodologie de Kimball repose sur l'organisation des données autour
d'une ou plusieurs \emph{tables de faits} entourées de \emph{dimensions}
descriptives.
\begin{itemize}
  \item Le \emph{schéma en étoile} centre une table de faits unique
        autour de dimensions dénormalisées.
  \item Le \emph{schéma en flocon} normalise les hiérarchies des
        dimensions (économie d'espace, mais jointures plus complexes).
  \item La \emph{constellation} (\emph{galaxy schema}) intègre plusieurs
        tables de faits partageant des dimensions conformes.
\end{itemize}

\subsection{Faits et grain}
Chaque table de faits est définie par son \emph{grain}~: le niveau de
détail le plus fin de chaque ligne. Un grain mal choisi conduit à des
agrégations erronées. Les mesures peuvent être~:
\begin{description}[leftmargin=2cm]
  \item[Additives.] Sommables sur toutes les dimensions
        (ex.\ \texttt{net\_amount}).
  \item[Semi-additives.] Sommables sur certaines dimensions seulement
        (ex.\ stock cumulé sur le temps).
  \item[Non-additives.] Ratios, pourcentages~; nécessitent un calcul
        spécifique (ex.\ \texttt{discount\_rate}).
\end{description}

\subsection{Dimensions à variation lente (SCD)}
Les attributs descriptifs évoluent dans le temps. Kimball identifie
plusieurs types de \ac{SCD}~:
\begin{itemize}
  \item \textbf{Type~1}~: écrasement (perte de l'historique).
  \item \textbf{Type~2}~: nouvelle ligne avec \texttt{valid\_from},
        \texttt{valid\_to}, \texttt{is\_current} (historisation
        complète).
  \item \textbf{Type~3}~: deux colonnes (valeur courante / précédente).
\end{itemize}
Notre système applique systématiquement le \textbf{\ac{SCD} Type~2} sur
les dimensions hors \texttt{dim\_date}.

\subsection{Surrogate keys}
Les \emph{surrogate keys} (suffixées \texttt{\_sk}) sont des clés
techniques générées par le \ac{DW}, indépendantes des clés métier
\textit{(natural keys)} issues des sources. Elles permettent d'absorber
les changements et de maintenir l'intégrité référentielle.

\section{Pipelines ETL et qualité des données}

\subsection{Phases de l'ETL}
Un pipeline \ac{ETL} comporte classiquement~:
\begin{enumerate}
  \item \textbf{Extraction}~: lecture depuis les sources hétérogènes.
  \item \textbf{Transformation}~: nettoyage, conversion, agrégation,
        gestion des doublons et des nulls.
  \item \textbf{Chargement}~: insertion dans les tables cibles avec
        gestion des \emph{upserts} et de la \ac{SCD}.
  \item \textbf{Réconciliation}~: contrôle des volumétries et des
        sommes de contrôle.
\end{enumerate}

\subsection{Change Data Capture (CDC)}
Le \ac{CDC} permet de ne charger que les changements depuis le dernier
batch via des \emph{watermarks} (timestamps ou colonnes monotones).
Cela divise par plusieurs ordres de grandeur la fenêtre de
chargement.

\subsection{Indicateurs de qualité}
Le profilage des données fournit, par colonne~: le pourcentage de nulls,
le nombre de valeurs uniques, la distribution des longueurs et les
valeurs extrêmes. Un \emph{score \ac{DQ}} agrégé permet de bloquer
le pipeline en cas d'anomalie majeure.

\section{Systèmes multi-agents et orchestration}

\subsection{Définition d'un agent IA}
Un \emph{agent IA} est une entité logicielle autonome dotée d'une boucle
\emph{perception--décision--action}. Dans le contexte des \ac{LLM}, un
agent étend les capacités d'un modèle conversationnel par~:
\begin{itemize}
  \item un état persistant (\emph{memory})~;
  \item un accès à des outils externes (\emph{tool calling})~;
  \item une logique de routage entre étapes (\emph{controller}).
\end{itemize}

\subsection{LangGraph et les graphes d'états}
\texttt{LangGraph} est un framework Python qui modélise un workflow
agentique comme un graphe orienté à arêtes conditionnelles. Chaque
nœud encapsule une fonction Python qui lit et met à jour un \emph{state}
partagé. Les fonctionnalités clés sont~:
\begin{itemize}
  \item \emph{checkpointing} (reprise après crash)~;
  \item \emph{interruption} avant un nœud
        (\texttt{interrupt\_before}~: utilisée pour notre étape
        \ac{HITL})~;
  \item \emph{stream mode} (diffusion incrémentale des mises à jour)~;
  \item \emph{cycles} (boucles de critique~/ correction).
\end{itemize}

\subsection{Human-in-the-Loop}
Le pattern \ac{HITL} maintient un humain comme arbitre final aux
points décisifs. Dans notre système, le pipeline s'interrompt après
la modélisation~: l'utilisateur peut soit valider, soit demander une
modification en langage naturel. Cette interaction est cruciale pour
la confiance dans le système et la conformité aux exigences métier.

\section{Modèles de langage et génération de code}

\subsection{Architectures et capacités}
Les \ac{LLM} modernes (GPT-4, Claude, GLM-5) reposent sur l'architecture
\emph{transformer} et excellent dans la compréhension du langage
naturel et la génération de code structuré. Pour notre cas d'usage,
nous exploitons trois capacités~:
\begin{itemize}
  \item compréhension de schémas relationnels en langage naturel~;
  \item génération de \ac{DDL} \ac{TSQL} cohérent~;
  \item production d'opérations atomiques au format \ac{JSON}
        (\emph{patch operations}) à partir d'une demande informelle.
\end{itemize}

\subsection{Cascade multi-fournisseurs}
Plutôt que de dépendre d'un fournisseur unique, notre système met en
œuvre une cascade~: GLM-5 (ZhipuAI), OpenRouter, Blaze, Ollama local,
puis un \emph{smart parser} déterministe en dernier recours. Cette
redondance garantit la continuité de service.

\subsection{Prompt engineering et patch operations}
La technique des \emph{patch operations} consiste à demander au \ac{LLM},
au lieu de régénérer un objet \ac{JSON} entier, de produire une liste
d'opérations atomiques (\texttt{add\_column}, \texttt{drop\_column},
\texttt{rename\_column}, \texttt{split\_date\_key},
\texttt{change\_column\_type}, \texttt{add\_fk}). Cette approche réduit
drastiquement la fenêtre de tokens et la probabilité d'erreur de
sérialisation.

\section{Conclusion}
Ce chapitre a posé les fondations conceptuelles du projet~: la
méthodologie Kimball pour la modélisation, les principes de l'\ac{ETL},
les architectures multi-agents et le rôle des \ac{LLM} dans la
génération de code. Le chapitre suivant présente les approches
existantes et positionne notre contribution dans l'état de l'art.

\cleardoublepage

% ============================================================
% CHAPITRE II : ETAT DE L'ART
% ============================================================
\chapter{État de l'art sur l'automatisation des entrepôts de données}
\label{chap:etat}

\section{Introduction}
Ce chapitre dresse un panorama des outils, frameworks et travaux de
recherche relatifs à la conception et à l'industrialisation
d'entrepôts de données. Nous analysons d'abord les outils
\emph{commerciaux} et \emph{open-source} dominants, puis nous présentons
les avancées récentes en agents \ac{LLM} pour l'ingénierie des données,
avant de positionner notre contribution.

\section{Plateformes ETL/ELT traditionnelles}

\subsection{dbt (data build tool)}
\texttt{dbt} est aujourd'hui la référence open-source pour la couche
\emph{transformation} d'un \ac{DW} cloud (Snowflake, BigQuery,
Redshift, Databricks). Il introduit une approche déclarative basée
sur des modèles \ac{SQL} versionnés en Git, des tests d'intégrité
intégrés et un graphe de dépendances calculé automatiquement.
Cependant, \texttt{dbt} n'automatise ni la modélisation initiale ni
le profilage~: l'utilisateur doit écrire chaque modèle à la main.

\subsection{Airbyte et Fivetran}
\texttt{Airbyte} (open-source) et \texttt{Fivetran} (commercial)
spécialisent l'\emph{extraction} et le \emph{chargement} (\ac{EL})~:
ils gèrent un large catalogue de connecteurs (300+) avec
synchronisation incrémentale. La transformation reste à la charge
de l'utilisateur, généralement via \texttt{dbt} en aval. Aucun de
ces outils ne propose de modélisation Kimball automatique.

\subsection{Talend et Pentaho Data Integration}
\texttt{Talend Open Studio} et \texttt{Pentaho} (\ac{ETL} graphique
historique) offrent une interface drag-and-drop pour composer des
pipelines complexes. Leur force est l'écosystème de transformations
prêtes à l'emploi~; leur faiblesse est la maintenance, la
collaboration Git et l'absence d'intelligence sémantique.

\subsection{Microsoft SSIS et Azure Data Factory}
Dans l'écosystème Microsoft, \texttt{SSIS} (\emph{SQL Server
Integration Services}) reste largement déployé pour les pipelines
\emph{on-premises}, tandis que \texttt{Azure Data Factory} est sa
version managée dans le cloud. Ces outils nécessitent une expertise
spécifique et n'intègrent aucune assistance par IA.

\subsection{Synthèse comparative}
Le tableau~\ref{tab:etl-tools} synthétise les caractéristiques
principales de ces plateformes au regard de cinq critères~: ingestion
multi-source, modélisation automatique, ETL automatique, validation
\ac{HITL} et déploiement \emph{cloud-native}.

\begin{table}[H]
\centering
\caption{Comparaison des principales plateformes ETL/ELT}
\label{tab:etl-tools}
\small
\begin{tabular}{lccccc}
\toprule
\textbf{Outil} & \textbf{Ingestion} & \textbf{Modélisation} & \textbf{ETL auto} & \textbf{HITL} & \textbf{Cloud} \\
\midrule
dbt           & --   & --   & ++   & --   & ++ \\
Airbyte       & ++   & --   & --   & --   & ++ \\
Fivetran      & ++   & --   & --   & --   & ++ \\
Talend        & ++   & --   & +    & +    & +  \\
SSIS / ADF    & ++   & --   & +    & --   & + (ADF) \\
\textbf{Notre système} & ++   & ++   & ++   & ++   & ++ \\
\bottomrule
\end{tabular}
\end{table}

\section{Apports de l'intelligence artificielle à l'ETL}

\subsection{Profilage automatique et détection d'anomalies}
Des travaux récents (\textit{e.g.}\ Great Expectations, Soda Core)
proposent des frameworks de validation déclarative pour la
qualité des données. L'intégration d'\ac{LLM} permet désormais
d'inférer automatiquement les attentes (\emph{expectations}) à
partir d'un échantillon, comme le démontre le projet Mostly~AI.

\subsection{Modélisation dimensionnelle assistée par LLM}
Plusieurs preuves de concept ont émergé en 2024-2025 pour automatiser
la modélisation Kimball~:
\begin{itemize}
  \item \texttt{Holistics AML}~: langage de modélisation déclaratif
        compilé en \ac{SQL}.
  \item \texttt{dbt Mesh + Cortex}~: extension de \texttt{dbt} avec
        suggestions de \ac{LLM} pour la création de modèles.
  \item \texttt{LangChain SQL Agent}~: agent générique capable de
        produire du \ac{SQL} en réponse à une question, sans
        toutefois construire un \ac{DW} complet.
\end{itemize}
Aucune de ces solutions ne couvre l'intégralité du cycle de vie
(profilage~$\to$~modélisation~$\to$~\ac{DDL}~$\to$~ETL~$\to$~lineage).

\subsection{Agents IA pour l'ingénierie des données}
Le paradigme \emph{agentic data engineering} émerge dans les
publications de 2024-2025 (Snowflake Data Cloud Summit, dbt Coalesce).
Les agents typiques sont~: agent d'ingestion, agent de transformation,
agent de qualité, agent de documentation. Toutefois, peu de projets
publient une implémentation complète orchestrée par graphe d'états
avec interruption \ac{HITL}.

\section{Validation humaine en boucle (HITL)}

\subsection{Patterns d'interaction}
Les patterns \ac{HITL} peuvent être classés selon le moment et la
profondeur de l'intervention~:
\begin{itemize}
  \item \textbf{Approbation}~: l'utilisateur dit oui/non.
  \item \textbf{Modification}~: l'utilisateur réécrit la sortie.
  \item \textbf{Conversation}~: l'utilisateur dialogue avec l'agent
        en langage naturel pour ajuster.
\end{itemize}
Notre système combine les trois~: validation binaire, modification
explicite via texte libre, et conversation Atlas pour les questions.

\subsection{Travaux notables}
\texttt{LangGraph Studio}, \texttt{AutoGen Studio} (Microsoft),
\texttt{CrewAI} fournissent des frameworks pour construire des
agents avec \ac{HITL}. Notre choix s'est porté sur \texttt{LangGraph}
pour sa maturité, sa gestion native du \emph{checkpointing} et sa
sémantique d'interruption précise (\texttt{interrupt\_before}).

\section{Reporting décisionnel et OLAP}

\subsection{Cubes OLAP et requêtage multidimensionnel}
Le requêtage \ac{OLAP} permet d'agréger les mesures selon plusieurs
dimensions (temps, produit, géographie). Les opérations classiques
sont \emph{slice}, \emph{dice}, \emph{drill-down} et \emph{roll-up}.
Notre système expose un \emph{cube logique} construit dynamiquement
à partir du modèle Kimball.

\subsection{Génération de rapports automatisée}
Des outils comme \texttt{Tableau Pulse} ou \texttt{Power BI Copilot}
introduisent la génération automatique d'insights et de visualisations.
Notre approche produit un fichier \texttt{.xlsx} natif à dix feuilles
incluant des graphiques embarqués (\texttt{openpyxl}), exploitable hors
ligne et sans licence.

\section{Positionnement de notre contribution}
Au regard de l'état de l'art, notre contribution se distingue sur
quatre axes~:
\begin{enumerate}
  \item \textbf{Automatisation \emph{end-to-end}}~: de la sauvegarde
        \texttt{.bak} brute jusqu'au rapport Excel décisionnel,
        sans intervention manuelle obligatoire.
  \item \textbf{Patch operations}~: alternative robuste à la
        régénération de schéma~JSON par les \ac{LLM}.
  \item \textbf{Smart parser déterministe}~: garantie de fonctionnement
        même sans \ac{LLM}, ce qui est unique dans la littérature.
  \item \textbf{Production-ready}~: déploiement Docker complet avec
        \ac{HTTPS} automatique, observabilité Prometheus/Grafana
        et 146 tests automatisés (unitaires, intégration, système,
        E2E).
\end{enumerate}

\section{Conclusion}
Ce chapitre a montré que l'écosystème actuel propose des outils puissants
mais fragmentés, et que les approches récentes basées sur les
\ac{LLM} restent embryonnaires. Notre projet comble cette lacune en
proposant une plateforme intégrée. Le chapitre suivant détaille la
conception et l'implémentation du système.

\cleardoublepage

% ============================================================
% CHAPITRE III : CONCEPTION ET IMPLEMENTATION
% ============================================================
\chapter{Conception et implémentation du système Agent Data Warehouse}
\label{chap:conception}

\section{Introduction}
Ce chapitre présente la conception détaillée et l'implémentation du
système. Nous décrivons d'abord l'architecture globale, puis les
agents spécialisés, l'assistant conversationnel Atlas, le \emph{smart
parser}, l'infrastructure de déploiement et enfin les résultats
expérimentaux obtenus sur la base \textit{Northwind}.

\section{Architecture globale}

\subsection{Vue macroscopique}
Le système est organisé selon une architecture trois-tiers~:
\begin{itemize}
  \item \textbf{Frontend}~: \emph{Single Page Application} React/Vite
        avec assistant flottant Atlas, panneau de validation HITL
        et grille d'exports.
  \item \textbf{Backend}~: API FastAPI exposant les endpoints REST
        et un flux \ac{SSE} temps réel, orchestré par
        \texttt{LangGraph}.
  \item \textbf{Data}~: SQL Server 2022 pour les métadonnées de
        session et l'entrepôt cible, plus un cache LLM en mémoire.
\end{itemize}

\subsection{Choix technologiques}
Le tableau~\ref{tab:techstack} récapitule la pile technologique.

\begin{table}[H]
\centering
\caption{Pile technologique du projet}
\label{tab:techstack}
\begin{tabular}{ll}
\toprule
\textbf{Couche} & \textbf{Technologies} \\
\midrule
Frontend  & React 18, Vite, Zustand, Tailwind, framer-motion \\
Backend   & FastAPI, Pydantic, Gunicorn + Uvicorn, SQLAlchemy \\
Agents IA & LangGraph 0.2, LangChain Core, requests \\
LLM       & GLM-5 (ZhipuAI), OpenRouter, Blaze, Ollama \\
Stockage  & SQL Server 2022, openpyxl (Excel) \\
Observabilité & Prometheus, Grafana, JSON logs \\
Déploiement & Docker multi-stage, Caddy (HTTPS Let's Encrypt) \\
Tests      & Pytest, Vitest, Playwright \\
\bottomrule
\end{tabular}
\end{table}

% ── Figure architecture globale ─────────────────────────────────
\begin{figure}[H]
\centering
\begin{tikzpicture}[
  font=\small,
  box/.style={rectangle, rounded corners=4pt, minimum width=3cm,
              minimum height=0.7cm, text centered,
              draw=black!60, fill=#1!15, drop shadow},
  arr/.style={-{Stealth[length=5pt]}, thick, color=black!50},
  every node/.style={font=\small}
]
% -- Couche Frontend
\node[box=blue, minimum width=8cm, minimum height=1.1cm]
      (fe) at (0,0) {\textbf{Frontend} — React/Vite + Atlas + Panneau HITL};
% -- Couche Backend
\node[box=orange, minimum width=8cm, minimum height=1.1cm]
      (be) at (0,-2) {\textbf{Backend} — FastAPI + LangGraph (23 nœuds) + SSE};
% -- Couche Data
\node[box=green, minimum width=3.5cm, minimum height=1.1cm]
      (db) at (-2.5,-4) {\textbf{SQL Server 2022}\\Entrepôt + Métadonnées};
\node[box=purple, minimum width=3.5cm, minimum height=1.1cm]
      (cache) at (2.5,-4) {\textbf{Cache LLM}\\LRU + TTL en mémoire};
% -- LLM Providers
\node[box=red, minimum width=8cm, minimum height=0.9cm]
      (llm) at (0,-5.8) {\textbf{Fournisseurs LLM} — GLM-5 · OpenRouter · Blaze · Ollama};
% -- Flèches
\draw[arr] (fe) -- node[right, font=\tiny]{REST/SSE} (be);
\draw[arr] (be) -- (db);
\draw[arr] (be) -- (cache);
\draw[arr] (be) -- (llm);
% -- Observabilité
\node[box=gray, minimum width=3cm, minimum height=0.7cm]
      (obs) at (5.2,-2) {Prometheus\\Grafana};
\draw[arr, dashed] (be) -- (obs);
\end{tikzpicture}
\caption{Architecture trois-tiers du système Agent Data Warehouse}
\label{fig:archi}
\end{figure}

\section{Pipeline multi-agents}

\subsection{Vue d'ensemble du graphe}
Le pipeline est modélisé comme un graphe \texttt{LangGraph} de 23 nœuds.
La figure~\ref{fig:pipeline} illustre le flux principal,
qui suit les phases canoniques d'un projet \ac{DW}~:

\begin{enumerate}
  \item \texttt{explorer}~: extraction des métadonnées sources.
  \item \texttt{data\_quality}~: profilage et calcul du score \ac{DQ}.
  \item \texttt{drift\_detector}~: détection des changements de schéma.
  \item \texttt{modeler}~: génération du modèle Kimball.
  \item \texttt{governance}~: contrôle de gouvernance.
  \item \texttt{critic}~: critique automatique du modèle.
  \item \texttt{human\_review}~: \emph{interrupt\_before} pour
        validation utilisateur.
  \item \texttt{chat\_modifier}~: application des modifications
        demandées (boucle avec \texttt{critic}).
  \item \texttt{cdc\_watermark}~: choix du mode (\emph{full\_load} ou
        incrémental).
  \item \texttt{etl\_tsql\_generator}~: génération du \ac{TSQL}.
  \item \texttt{etl\_initializer}, \texttt{etl\_extractor},
        \texttt{etl\_transformer}, \texttt{etl\_loader}~: phases
        d'exécution ETL.
  \item \texttt{healer}~: auto-réparation en cas d'erreur ETL.
  \item \texttt{lineage\_tracker}~: construction du graphe de lignage.
  \item \texttt{query\_generator}, \texttt{insight\_generator}~:
        analyses post-chargement.
  \item \texttt{cataloger}, \texttt{airflow\_generator},
        \texttt{dbt\_generator}~: production des artefacts.
\end{enumerate}

% ── Figure pipeline LangGraph ────────────────────────────────────
\begin{figure}[H]
\centering
\begin{tikzpicture}[
  font=\scriptsize,
  node/.style={rectangle, rounded corners=3pt, minimum width=2.6cm,
               minimum height=0.55cm, text centered, draw=black!60,
               fill=#1!18, drop shadow},
  arr/.style={-{Stealth[length=4pt]}, thick},
  skip/.style={-{Stealth[length=4pt]}, thick, dashed, color=red!70}
]
% Colonne centrale
\node[node=blue]   (exp)  at (0,0)     {explorer};
\node[node=blue]   (dq)   at (0,-1.0)  {data\_quality};
\node[node=blue]   (dd)   at (0,-2.0)  {drift\_detector};
\node[node=orange] (mod)  at (0,-3.0)  {modeler};
\node[node=orange] (gov)  at (0,-4.0)  {governance};
\node[node=orange] (cri)  at (0,-5.0)  {critic};
\node[node=red]    (hr)   at (0,-6.0)  {human\_review [HITL]};
\node[node=purple] (cm)   at (0,-7.0)  {chat\_modifier};
\node[node=blue]   (cdc)  at (0,-8.0)  {cdc\_watermark};
\node[node=green]  (tsql) at (0,-9.0)  {etl\_tsql\_generator};
\node[node=green]  (etl)  at (0,-10.0) {etl\_init / extract / transform / load};
\node[node=red]    (heal) at (0,-11.0) {healer};
\node[node=teal]   (lin)  at (0,-12.0) {lineage\_tracker};
\node[node=teal]   (rep)  at (0,-13.0) {query\_gen · insight\_gen · cataloger};
% Flèches principales
\foreach \a/\b in {exp/dq, dq/dd, dd/mod, mod/gov, gov/cri, cri/hr,
                    hr/cdc, cdc/tsql, tsql/etl, etl/heal, heal/lin, lin/rep}
  \draw[arr] (\a) -- (\b);
% Boucle critique → chat_modifier → critic
\draw[arr] (cri.east) -- ++(1.5,0) |- node[right, pos=0.25,
     font=\tiny]{rejet} (cm.east);
\draw[arr] (cm.west) -- ++(-1.5,0) |- (cri.west);
% Boucle healer
\draw[arr] (heal.east) -- ++(0.8,0) |- node[right, pos=0.25,
     font=\tiny]{retry} (etl.east);
\end{tikzpicture}
\caption{Graphe LangGraph du pipeline Agent Data Warehouse (23 nœuds)}
\label{fig:pipeline}
\end{figure}

\subsection{Routage conditionnel}
Le routage entre nœuds est implémenté par des fonctions
\texttt{route\_after\_*} qui inspectent l'état partagé. Par exemple,
\texttt{route\_after\_critic} renvoie soit \texttt{human\_review}
(modèle approuvé), soit \texttt{chat\_modifier} (modèle rejeté), avec
un compteur de boucle pour éviter les itérations infinies.

\subsection{Gestion des erreurs et auto-réparation}
En cas d'échec d'une étape ETL (insertion contrainte, type incompatible),
le routeur dirige vers le nœud \texttt{healer} qui analyse la stack
trace et tente jusqu'à \texttt{MAX\_RETRIES = 3} corrections. Au-delà,
le pipeline termine en erreur explicite.

\section{L'assistant conversationnel Atlas}

\subsection{Mode hybride conversation/modification}
Atlas détecte automatiquement l'intention de l'utilisateur (question
libre vs.\ modification du modèle) grâce à l'analyse de mots-clés et
de la structure de phrase. Les conversations libres reçoivent une
réponse type ChatGPT~; les demandes de modification déclenchent le
nœud \texttt{chat\_modifier}.

\subsection{Patch operations}
Au lieu de régénérer le modèle entier, Atlas demande au \ac{LLM}
de produire une liste d'opérations atomiques au format \ac{JSON}.
Le code~\ref{lst:patch} illustre une réponse type pour la demande
``ajoute net\_amount et \emph{role-playing} dates''.

\begin{lstlisting}[language=Python, caption={Exemple de patch operations produit par Atlas}, label=lst:patch]
{
  "ops": [
    {
      "op": "split_date_key",
      "table": "fact_orders",
      "old_column": "date_sk",
      "new_columns": [
        {"name": "order_date_sk", "nullable": false},
        {"name": "required_date_sk", "nullable": false},
        {"name": "shipped_date_sk", "nullable": true}
      ]
    },
    {
      "op": "add_column",
      "table": "fact_orders",
      "column": {
        "name": "net_amount",
        "type": "DECIMAL(15,4)",
        "role": "computed",
        "description": "unitprice * quantity * (1 - discount)"
      }
    },
    {
      "op": "change_column_type",
      "table": "dim_employee",
      "column": "reportsto",
      "type": "INT"
    }
  ],
  "summary": "Role-playing dates + net_amount + reportsto INT"
}
\end{lstlisting}

\subsection{Cascade multi-fournisseurs}
La fonction \texttt{get\_human\_review\_llm()} tente successivement~:
\begin{enumerate}
  \item GLM-5 (ZhipuAI) avec rotation sur~3 clés~;
  \item OpenRouter avec un modèle gratuit ou premium~;
  \item Blaze (si configuré)~;
  \item Ollama local avec modèles légers (\texttt{phi3:mini},
        \texttt{qwen2.5:0.5b})~;
  \item enfin, le \emph{smart parser} déterministe.
\end{enumerate}
Cette redondance garantit que la modification est appliquée même en
cas de panne d'un fournisseur \ac{LLM}.

\section{Smart Parser : modification sans LLM}
\label{sec:smartparser}

\subsection{Principe}
Le \emph{smart parser} est un module Python pur qui parse une demande
multi-étapes (numérotée 1./2./3.) et infère des opérations atomiques
sans recourir à un \ac{LLM}. Il s'appuie sur des heuristiques
métier~:

\begin{itemize}
  \item Si la demande contient ``multi-dates'' ou ``role-playing''
        avec trois noms en \texttt{*\_date\_sk}, le parser produit
        une opération \texttt{split\_date\_key}.
  \item Si la demande mentionne ``mesure calculée'' avec un nom
        capitalisé entre crochets, le parser produit une opération
        \texttt{add\_column} de rôle \texttt{computed}.
  \item Si la demande évoque un changement de type
        (``change le type de X en INT''), une opération
        \texttt{change\_column\_type} est émise.
  \item Les ``contraintes de clés étrangères'' déclenchent une série
        d'opérations \texttt{add\_fk}.
\end{itemize}

\subsection{Évaluation expérimentale}
Sur la demande Kimball complexe de quatre directives utilisée comme
test, le \emph{smart parser} infère correctement six opérations~:
\begin{enumerate}
  \item split de \texttt{date\_sk} en trois clés (avec
        \texttt{nullable=true} sur \texttt{shipped\_date\_sk})~;
  \item ajout de \texttt{net\_amount DECIMAL(15,4)}~;
  \item changement de type de \texttt{reportsto} en \texttt{INT}~;
  \item trois \texttt{add\_fk} entre la table de faits et les
        dimensions \texttt{dim\_product}, \texttt{dim\_customer},
        \texttt{dim\_shipper}.
\end{enumerate}
Cette robustesse permet de garantir le bon fonctionnement du système
indépendamment de l'état des fournisseurs \ac{LLM}.

\section{Génération du DDL T-SQL}
La fonction \texttt{\_generate\_ddl()} traduit le modèle logique \ac{JSON}
en \ac{TSQL} exécutable sur SQL Server~2022. Elle prend en charge~:
\begin{itemize}
  \item les \emph{surrogate keys} (\texttt{BIGINT IDENTITY(1,1)
        PRIMARY KEY})~;
  \item les colonnes de \ac{SCD} Type~2 (\texttt{valid\_from},
        \texttt{valid\_to}, \texttt{is\_current}, \texttt{row\_hash})~;
  \item les index uniques filtrés sur \texttt{is\_current = 1}~;
  \item les index \texttt{NONCLUSTERED} sur les \ac{FK} de la table
        de faits~;
  \item les tables de quarantaine pour les rejets ETL.
\end{itemize}

\section{Reporting décisionnel}

\subsection{Rapport Excel à dix feuilles}
Le module \texttt{export\_service.py} génère un fichier Excel décisionnel
avec dix feuilles~: tableau de bord, mesures \& KPI, qualité données,
schéma étoile, performance ETL, analyses OLAP, catalogue, lignage,
DDL, journal d'exécution. La feuille \emph{Mesures \& KPI} embarque
des graphiques natifs (bar, line, pie) calculés sur des agrégations
multi-granularité (année, trimestre, mois) via \texttt{openpyxl.chart}.

\subsection{Backup .bak et fallback logique}
L'export \texttt{.bak} déclenche un \texttt{BACKUP DATABASE} natif
SQL Server. En cas d'indisponibilité, un \emph{backup logique}
(archive \texttt{.zip} contenant le \ac{DDL} et un export \ac{CSV}
de chaque table) est généré comme \emph{fallback}, garantissant
que l'utilisateur reçoit toujours un livrable.

\section{Infrastructure de déploiement}

\subsection{Docker multi-stage}
Le \texttt{Dockerfile.backend} utilise une approche multi-étapes
(\emph{builder} + runtime slim) qui produit une image finale
$\sim$~600~MB. Les étapes clés sont~:
\begin{itemize}
  \item \emph{builder}~: installation des dépendances de compilation
        et création d'un \texttt{venv} isolé~;
  \item \emph{runtime}~: image \texttt{python:3.11-slim} avec uniquement
        les dépendances système nécessaires
        (\texttt{msodbcsql18}, \texttt{tini})~;
  \item \emph{user}~: utilisateur non-root \texttt{appuser} (UID 1000)~;
  \item \emph{server}~: \texttt{gunicorn} avec workers \texttt{uvicorn}.
\end{itemize}

\subsection{Reverse-proxy Caddy avec HTTPS automatique}
Le service \texttt{caddy} déployé en frontal obtient un certificat
Let's Encrypt automatiquement via le protocole ACME, supporte
HTTP/3 et applique les en-têtes de sécurité (HSTS,
\texttt{Strict-Transport-Security}, etc.).

\subsection{Observabilité}
Un endpoint \texttt{/metrics} expose les compteurs et histogrammes
Prometheus (requêtes par route, latence P50/P95/P99, durée des
agents). Un dashboard Grafana JSON pré-provisionné visualise ces
métriques. Les logs peuvent être émis au format \ac{JSON} structuré
pour ingestion ELK/Loki.

\section{Tests automatisés}
La couverture de tests est répartie sur cinq niveaux~:
\begin{itemize}
  \item \textbf{Unitaires backend}~: 94 tests \texttt{pytest} sur
        \texttt{chat\_modifier}, \texttt{modeler}, \texttt{llm\_factory}
        et \texttt{export\_service}.
  \item \textbf{Unitaires frontend}~: 29 tests \texttt{vitest}
        sur le store Zustand et les composants React.
  \item \textbf{Intégration}~: 27 tests sur les endpoints REST
        (validation, sécurité, multi-tenant, email).
  \item \textbf{Système}~: 9 tests sur le flux complet du pipeline.
  \item \textbf{End-to-end}~: 13 tests \texttt{Playwright} sur l'UI
        réelle (ouverture Atlas, suggestions, exports).
\end{itemize}
Le total atteint \textbf{146 tests~Python} et \textbf{42 tests~JS},
exécutés en CI GitHub Actions à chaque \emph{push}.

\section{Sécurité}
Le système intègre plusieurs mécanismes de sécurité~:
\begin{itemize}
  \item authentification \ac{JWT} avec rotation possible du secret~;
  \item rate-limiting via \texttt{slowapi}~;
  \item en-têtes de sécurité (X-Frame-Options, X-Content-Type-Options,
        Permissions-Policy, CSP)~;
  \item validation \texttt{Pydantic} stricte des entrées~;
  \item garde anti-injection \ac{SQL} sur l'endpoint
        \texttt{/api/execute-query} (refus de tout sauf
        \texttt{SELECT})~;
  \item isolation multi-tenant par \texttt{user\_prefix}.
\end{itemize}

\section{Résultats expérimentaux}

\subsection{Cas d'étude : Northwind}
Nous avons exécuté le pipeline complet sur une sauvegarde \texttt{.bak}
de la base de référence \textit{Northwind} (Microsoft). Le
schéma en étoile généré est représenté à la figure~\ref{fig:star}.
Les résultats obtenus sont synthétisés dans le tableau~\ref{tab:results}.

% ── Figure schéma en étoile Northwind ───────────────────────────
\begin{figure}[H]
\centering
\begin{tikzpicture}[
  font=\scriptsize,
  fact/.style={rectangle, rounded corners=3pt, minimum width=3.8cm,
               text width=3.6cm, align=center, minimum height=2.2cm,
               draw=orange!80!black, fill=orange!20, drop shadow},
  dim/.style={rectangle, rounded corners=3pt, minimum width=3cm,
              text width=2.8cm, align=center, minimum height=1.5cm,
              draw=blue!70, fill=blue!10, drop shadow},
  fk/.style={-{Stealth[length=4pt]}, thick, color=black!50}
]
% Table de faits (centre)
\node[fact] (fo) at (0,0) {
  \textbf{fact\_orders}\\[3pt]
  order\_sk (PK)\\order\_date\_sk (FK)\\required\_date\_sk (FK)\\shipped\_date\_sk (FK)\\customer\_sk (FK)\\employee\_sk (FK)\\product\_sk (FK)\\shipper\_sk (FK)\\quantity · unit\_price\\discount · net\_amount
};
% Dimensions
\node[dim] (date) at (0, 4.5)  {\textbf{dim\_date}\\date\_sk (PK)\\date · year\\quarter · month\\day · weekday};
\node[dim] (cust) at (-5.5, 2) {\textbf{dim\_customer}\\customer\_sk (PK)\\customer\_id\\company\_name\\country · city};
\node[dim] (emp)  at (-5.5,-2) {\textbf{dim\_employee}\\employee\_sk (PK)\\full\_name\\title · region\\valid\_from/to};
\node[dim] (prod) at (5.5, 2)  {\textbf{dim\_product}\\product\_sk (PK)\\product\_name\\category · supplier\\unit\_price};
\node[dim] (ship) at (5.5,-2)  {\textbf{dim\_shipper}\\shipper\_sk (PK)\\company\_name\\phone};
% Flèches FK
\draw[fk] (fo.north)      -- (date.south);
\draw[fk] (fo.west)       -- (cust.east);
\draw[fk] (fo.south west) -- (emp.east);
\draw[fk] (fo.east)       -- (prod.west);
\draw[fk] (fo.south east) -- (ship.west);
\end{tikzpicture}
\caption{Schéma en étoile généré pour la base \emph{Northwind}}
\label{fig:star}
\end{figure}

\begin{table}[H]
\centering
\caption{Résultats du pipeline sur la base \emph{Northwind}}
\label{tab:results}
\begin{tabular}{lr}
\toprule
\textbf{Indicateur} & \textbf{Valeur} \\
\midrule
Tables sources                          & 14 \\
Score qualité (DQ)                      & 97/100 \\
Tables de faits générées                & 1 (\texttt{fact\_orders}) \\
Dimensions générées                     & 5 \\
Lignes de DDL produites                 & $\approx$\ 250 \\
Caractères de DDL                       & 8 037 \\
Tests passants                          & 146/146 (100~\%) \\
Durée pipeline complet (sans LLM)       & $\approx$\ 6 min \\
Taille image Docker backend             & $\approx$\ 600 MB \\
\bottomrule
\end{tabular}
\end{table}

\subsection{Modifications appliquées via Atlas}
La directive Kimball à quatre points (multi-dates, mesure calculée
\texttt{net\_amount}, type de \texttt{reportsto}, contraintes \ac{FK})
a été correctement traduite en six opérations atomiques par le
\emph{smart parser} et appliquée au modèle, démontrant l'efficacité
de l'approche \emph{patch operations} décrite en
section~\ref{sec:smartparser}.

\section{Limites et travaux futurs}

\subsection{Limites identifiées}
\begin{itemize}
  \item Le système est aujourd'hui couplé à SQL Server~; un support
        multi-base (Postgres, Snowflake, BigQuery) nécessiterait
        une couche d'abstraction supplémentaire.
  \item Le \emph{smart parser} est déterministe sur un sous-ensemble
        d'opérations courantes~; les demandes très spécifiques
        nécessitent encore un \ac{LLM}.
  \item Les performances ETL sur de gros volumes (>~100~M de lignes)
        n'ont pas été éprouvées.
\end{itemize}

\subsection{Pistes d'amélioration}
\begin{itemize}
  \item Connecteurs natifs Snowflake et BigQuery.
  \item Mode collaboratif multi-utilisateurs avec notifications
        temps réel.
  \item Intégration Power~BI \emph{push API} pour publication
        automatique du dataset.
  \item Plug-in VS Code pour interagir avec Atlas depuis l'IDE.
\end{itemize}

\section{Conclusion}
Ce chapitre a détaillé l'architecture, les composants et les résultats
du système \emph{Agent Data Warehouse}. La combinaison du paradigme
multi-agents, des \emph{patch operations}, du \emph{smart parser}
et d'une infrastructure \emph{production-ready} constitue, à notre
connaissance, la première implémentation complète et open-source de
ce type.

\cleardoublepage

% ────────────────────────────────────────────────────────────
% CONCLUSION GENERALE
% ────────────────────────────────────────────────────────────
\chapter*{Conclusion générale}
\addcontentsline{toc}{chapter}{Conclusion générale}
\markboth{Conclusion générale}{}

Ce mini-projet a abordé la problématique de l'industrialisation des
entrepôts de données à travers le prisme de l'intelligence artificielle
agentique. La complexité, le coût et les délais inhérents à un projet
\ac{DW} traditionnel constituent un frein notable à l'adoption de la
\ac{BI} par de nombreuses organisations. Nous avons proposé une
plateforme intégrée, \emph{Agent Data Warehouse}, qui automatise
l'intégralité du cycle de vie tout en préservant le contrôle humain
aux étapes critiques.

\paragraph{Synthèse des contributions.}
Les apports principaux de ce travail sont au nombre de cinq~:
\begin{enumerate}
  \item une architecture \emph{multi-agents} de 23 nœuds spécialisés
        orchestrés par \texttt{LangGraph}, couvrant l'ingestion,
        la modélisation, l'exécution \ac{ETL}, le lignage et le
        reporting~;
  \item un assistant conversationnel \textbf{Atlas} fondé sur le
        paradigme \emph{patch operations}, permettant des
        modifications du schéma en langage naturel avec une
        cascade robuste de fournisseurs \ac{LLM}~;
  \item un \emph{smart parser} déterministe garantissant la
        continuité de service même en cas de défaillance des
        \ac{LLM}, capable d'inférer six opérations atomiques à
        partir d'une demande Kimball complexe~;
  \item un rapport décisionnel Excel à dix feuilles avec
        agrégations multi-granularité et graphiques natifs~;
  \item une infrastructure \emph{production-ready} (Docker
        multi-étapes, \ac{HTTPS} automatique via Caddy,
        observabilité Prometheus/Grafana) et une couverture de
        tests automatisés (146 tests Python).
\end{enumerate}

\paragraph{Validation expérimentale.}
L'expérimentation sur la base de référence \textit{Northwind} a
démontré la capacité du système à produire, sans intervention
manuelle, un schéma en étoile complet (1~table de faits +
5~dimensions), un \ac{DDL} \ac{TSQL} de 8\,037~caractères, un score
qualité \ac{DQ} de 97/100 et un rapport Excel décisionnel
exploitable. La directive Kimball à quatre points formulée en
langage naturel a été correctement traduite en six opérations
atomiques appliquées au modèle.

\paragraph{Perspectives.}
Plusieurs axes d'évolution se dégagent~: (i) extension à des cibles
cloud (Snowflake, BigQuery) via une couche d'abstraction~;
(ii) optimisation pour des volumes massifs (partitioning, index
\emph{columnstore})~; (iii) intégration directe à des outils \ac{BI}
(Power~BI, Tableau) via leurs API \emph{push}~; (iv) extension du
\emph{smart parser} pour couvrir des cas d'usage plus avancés
(constellations, dimensions junk, \ac{SCD}~Type~3)~;
(v) instrumentation d'une métrique de \emph{coût total de
possession} (TCO) comparant notre approche à un projet \ac{DW}
classique sur des cas d'études réels.

\paragraph{Conclusion personnelle.}
Ce projet a été pour moi l'occasion de mettre en pratique les
connaissances acquises durant ma formation en \ac{SIAD} et
d'explorer des technologies de pointe (LangGraph, modèles
de langage, conteneurisation, observabilité). Au-delà des aspects
techniques, j'ai pu mesurer l'importance du couplage entre
intelligence artificielle et expertise humaine~: la valeur d'un
système ne réside pas seulement dans son automatisation, mais aussi
dans sa capacité à s'effacer devant le jugement de l'utilisateur
aux moments critiques. Cette philosophie de l'\emph{Human-in-the-Loop}
me semble être la voie la plus prometteuse pour l'avenir de
l'ingénierie des données.

\cleardoublepage

% ────────────────────────────────────────────────────────────
% BIBLIOGRAPHIE
% ────────────────────────────────────────────────────────────
\addcontentsline{toc}{chapter}{Bibliographie}
\begin{thebibliography}{99}

\bibitem{kimball2013}
Ralph Kimball, Margy Ross.
\textit{The Data Warehouse Toolkit: The Definitive Guide to
Dimensional Modeling}.
3\textsuperscript{rd} edition, Wiley, 2013.

\bibitem{inmon2005}
William H. Inmon.
\textit{Building the Data Warehouse}.
4\textsuperscript{th} edition, Wiley, 2005.

\bibitem{kimball2008lifecycle}
Ralph Kimball, Margy Ross, Warren Thornthwaite, Joy Mundy, Bob Becker.
\textit{The Data Warehouse Lifecycle Toolkit}.
2\textsuperscript{nd} edition, Wiley, 2008.

\bibitem{vassiliadis2009}
Panos Vassiliadis.
\textit{A Survey of Extract-Transform-Load Technology}.
International Journal of Data Warehousing and Mining, 5(3), 1-27, 2009.

\bibitem{langgraph2024}
LangChain Inc.
\textit{LangGraph: Building Stateful, Multi-Actor LLM Applications}.
\url{https://langchain-ai.github.io/langgraph/}, 2024.

\bibitem{vaswani2017}
Ashish Vaswani et al.
\textit{Attention Is All You Need}.
Advances in Neural Information Processing Systems (NeurIPS), 2017.

\bibitem{brown2020}
Tom B. Brown et al.
\textit{Language Models are Few-Shot Learners}.
NeurIPS, 2020.

\bibitem{dbt2023}
Tristan Handy et al.
\textit{Analytics Engineering with dbt}.
\url{https://docs.getdbt.com/}, 2023.

\bibitem{microsoft_tsql}
Microsoft.
\textit{Transact-SQL Reference (Database Engine)}.
\url{https://learn.microsoft.com/en-us/sql/t-sql/}, 2024.

\bibitem{fastapi}
Sebastián Ramírez.
\textit{FastAPI Documentation}.
\url{https://fastapi.tiangolo.com/}, 2024.

\bibitem{caddy}
Matt Holt et al.
\textit{Caddy: The Ultimate Server with Automatic HTTPS}.
\url{https://caddyserver.com/}, 2024.

\bibitem{prometheus2023}
Prometheus Authors.
\textit{Prometheus: Monitoring System and Time Series Database}.
\url{https://prometheus.io/}, 2023.

\bibitem{grafana2024}
Grafana Labs.
\textit{Grafana: The open observability platform}.
\url{https://grafana.com/}, 2024.

\bibitem{ollama2024}
Ollama.
\textit{Ollama: Run large language models locally}.
\url{https://ollama.com/}, 2024.

\bibitem{zhipu_glm}
ZhipuAI.
\textit{GLM-4 Technical Report}.
\url{https://github.com/THUDM/GLM-4}, 2024.

\bibitem{openrouter}
OpenRouter.
\textit{OpenRouter API Documentation}.
\url{https://openrouter.ai/docs}, 2024.

\bibitem{wei2022cot}
Jason Wei et al.
\textit{Chain-of-Thought Prompting Elicits Reasoning in Large
Language Models}.
NeurIPS, 2022.

\bibitem{yao2023react}
Shunyu Yao et al.
\textit{ReAct: Synergizing Reasoning and Acting in Language Models}.
ICLR, 2023.

\bibitem{playwright}
Microsoft.
\textit{Playwright: End-to-end testing for modern web apps}.
\url{https://playwright.dev/}, 2024.

\bibitem{docker}
Docker Inc.
\textit{Docker Documentation}.
\url{https://docs.docker.com/}, 2024.

\bibitem{soda2024}
Soda Data.
\textit{Soda Core: Open-source data quality framework}.
\url{https://github.com/sodadata/soda-core}, 2024.

\bibitem{ge2024}
Great Expectations.
\textit{Great Expectations: Data Quality Tool}.
\url{https://greatexpectations.io/}, 2024.

\end{thebibliography}

\cleardoublepage

% ============================================================
% ANNEXES
% ============================================================
\appendix

\chapter{Guide d'installation et de démarrage rapide}
\label{ann:install}

\section{Prérequis}
\begin{itemize}
  \item Docker Desktop $\geq$ 24.0 et Docker Compose V2.
  \item SQL Server 2022 ou accès à une instance distante.
  \item Python 3.11+ (pour le développement local sans Docker).
  \item Node.js 20+ et npm 10+ (pour le frontend).
  \item Au moins une clé API valide~: GLM-5 (ZhipuAI) ou OpenRouter.
\end{itemize}

\section{Démarrage via Docker Compose}
\begin{lstlisting}[language=bash, caption={Démarrage de la stack complète}]
# 1. Cloner le dépôt
git clone https://github.com/votre-org/agent-dw.git
cd agent-dw

# 2. Configurer les variables d'environnement
cp .env.example .env
# Editer .env : renseigner ZHIPU_API_KEY, SQL_SERVER, etc.

# 3. Lancer la stack (backend + frontend + SQL Server + Caddy)
docker compose up -d --build

# 4. Vérifier l'état des conteneurs
docker compose ps

# 5. Accéder à l'interface
# --> https://localhost (HTTPS automatique via Caddy)
\end{lstlisting}

\section{Démarrage en mode développement local}
\begin{lstlisting}[language=bash, caption={Mode développement sans Docker}]
# Backend
python -m venv .venv && source .venv/Scripts/activate  # Windows
pip install -r requirements.txt
uvicorn main:app --reload --port 8000

# Frontend (dans un second terminal)
npm install
npm run dev
# --> http://localhost:5173
\end{lstlisting}

\section{Variables d'environnement clés}
\begin{table}[H]
\centering
\caption{Variables d'environnement principales}
\label{tab:env}
\begin{tabular}{lp{8cm}}
\toprule
\textbf{Variable} & \textbf{Description} \\
\midrule
\texttt{ZHIPU\_API\_KEY\_1..3} & Clés API GLM-5 (ZhipuAI) avec rotation automatique \\
\texttt{OPENROUTER\_API\_KEY} & Clé OpenRouter (fournisseur secondaire) \\
\texttt{SQL\_SERVER} & Hôte SQL Server (ex. \texttt{localhost,1433}) \\
\texttt{SQL\_DATABASE} & Nom de la base de staging/métadonnées \\
\texttt{SECRET\_KEY} & Clé JWT (256 bits aléatoires recommandés) \\
\texttt{MAX\_RETRIES} & Nombre max de tentatives du \emph{healer} (défaut~: 3) \\
\texttt{LLM\_CACHE\_TTL} & Durée de vie du cache LLM en secondes (défaut~: 3600) \\
\texttt{PROMETHEUS\_ENABLED} & Activer l'export des métriques (\texttt{true}/\texttt{false}) \\
\bottomrule
\end{tabular}
\end{table}

% ────────────────────────────────────────────────────────────────
\chapter{Extraits de code significatifs}
\label{ann:code}

\section{Structure du nœud \texttt{modeler}}
\begin{lstlisting}[language=Python,
  caption={Nœud modeler -- génération du modèle Kimball (extrait)},
  label=lst:modeler]
def modeler(state: DWState) -> DWState:
    """Agent de modélisation dimensionnelle (Kimball)."""
    schema_info = state["schema_info"]
    dq_report   = state["dq_report"]

    llm = get_modeler_llm()          # Cascade GLM-5 -> OpenRouter -> ...
    prompt = build_kimball_prompt(schema_info, dq_report)

    raw = llm.invoke(prompt)
    model = parse_kimball_model(raw)  # Validation Pydantic stricte

    return {
        **state,
        "kimball_model": model,
        "critic_cycles":  0,
        "status":        "modeled",
    }
\end{lstlisting}

\section{Fonction de cascade LLM}
\begin{lstlisting}[language=Python,
  caption={Cascade multi-fournisseurs LLM (extrait simplifié)},
  label=lst:cascade]
def get_human_review_llm(config: AppConfig):
    """Retourne le premier LLM disponible parmi la cascade."""
    providers = [
        lambda: _try_glm5(config),
        lambda: _try_openrouter(config),
        lambda: _try_blaze(config),
        lambda: _try_ollama(config),
    ]
    for provider_fn in providers:
        try:
            llm = provider_fn()
            llm.invoke("ping")   # Test de connectivite rapide
            return llm
        except Exception as exc:
            logger.warning("Provider failed: %s", exc)
    # Dernier recours : smart parser deterministe
    return SmartParserFallback()
\end{lstlisting}

\section{Extrait du \emph{Smart Parser}}
\begin{lstlisting}[language=Python,
  caption={Smart Parser -- detection de split de dates (extrait)},
  label=lst:smartparser]
import re

DATE_SPLIT_PATTERN = re.compile(
    r"(multi.?dates?|role.?playing)",
    re.IGNORECASE
)

def _parse_date_split(directive: str, table: str) -> list[dict]:
    """Infere une operation split_date_key depuis une directive texte."""
    if not DATE_SPLIT_PATTERN.search(directive):
        return []
    # Recherche des noms de colonnes *_date_sk
    cols = re.findall(r"(\w+_date_sk)", directive)
    if len(cols) < 2:
        return []
    return [{
        "op":         "split_date_key",
        "table":      table,
        "old_column": "date_sk",
        "new_columns": [
            {"name": c, "nullable": "shipped" in c}
            for c in cols
        ],
    }]
\end{lstlisting}

% ────────────────────────────────────────────────────────────────
\chapter{Configuration Docker Compose (extrait)}
\label{ann:docker}

\begin{lstlisting}[language=bash,
  caption={docker-compose.yml -- services principaux (extrait)},
  label=lst:docker]
services:

  backend:
    build:
      context: .
      dockerfile: docker/Dockerfile.backend
    environment:
      - SQL_SERVER=${SQL_SERVER}
      - ZHIPU_API_KEY_1=${ZHIPU_API_KEY_1}
      - SECRET_KEY=${SECRET_KEY}
    ports:
      - "8000:8000"
    depends_on:
      sqlserver:
        condition: service_healthy
    restart: unless-stopped

  frontend:
    build:
      context: .
      dockerfile: docker/Dockerfile.frontend
    ports:
      - "3000:80"
    restart: unless-stopped

  caddy:
    image: caddy:2-alpine
    ports:
      - "80:80"
      - "443:443"
    volumes:
      - ./docker/Caddyfile:/etc/caddy/Caddyfile:ro
      - caddy_data:/data
    restart: unless-stopped

  sqlserver:
    image: mcr.microsoft.com/mssql/server:2022-latest
    environment:
      - ACCEPT_EULA=Y
      - SA_PASSWORD=${SA_PASSWORD}
    healthcheck:
      test: ["/opt/mssql-tools/bin/sqlcmd",
             "-S", "localhost", "-U", "sa",
             "-P", "${SA_PASSWORD}", "-Q", "SELECT 1"]
      interval: 10s
      retries: 10
    restart: unless-stopped

volumes:
  caddy_data:
\end{lstlisting}

\end{document}
